\documentclass[12pt]{article}
\usepackage[utf8]{inputenc}
\usepackage[T1]{fontenc}
\usepackage{hyperref}
\title{Kaninuwa Language Resources for K-12 Curriculum Development}
\author{Compiled by Claude Sonnet 5 (claude-sonnet-5) for Andr\'as Kornai}
\date{}
\begin{document}
\maketitle

\section{Introduction}
Kaninuwa (also called Wataluma; ISO 639-3 \texttt{wat}, Glottocode kani1281) is an Oceanic language of northern Goodenough Island, Milne Bay, Papua New Guinea. Search results give roughly 300 to 360 speakers, based on older figures from around 2000. Very little documentation was found.

\section{Primary Educational Resources}
No textbooks or primers were found. Joshua Project reports vernacular education in the first three years of primary school, which I could not verify elsewhere:
\begin{itemize}
\item Joshua Project, Wataluma in Papua New Guinea: \url{https://joshuaproject.net/people_groups/15887}
\end{itemize}

\section{Linguistic and Reference Materials}
\begin{itemize}
\item Wikipedia, Kaninuwa language (overview and alphabet): \url{https://en.wikipedia.org/wiki/Kaninuwa_language}
\item Omniglot page on Kaninuwa: \url{https://omniglot.com/writing/kaninuwa.htm}
\item Glosbe Kaninuwa-to-Nimowa listing, a crowd-sourced word list of unverified quality: \url{https://glosbe.com/wat/nmx}
\end{itemize}

\section{Religious and Community Resources}
Joshua Project (above) states that there is no Bible in the mother tongue.

\section{Digital and Audio Resources}
No audio or digital resources were found.

\section{Curriculum Development Context}
The resources found are general web descriptions. No curriculum materials, dictionary or grammar were found. Any curriculum would require fieldwork and local partnership, and the Papua New Guinea education authorities and local mission or school contacts would be the next place to ask.

\end{document}